\documentclass[11 pt]{article}  % sets the font to 12 pt and says this is an 
                                % article (as opposed to book or other documents)
\usepackage{amsfonts, amssymb}  % math fonts and symbols
\usepackage{amsmath}            % math
\usepackage{enumitem}           % change enumeration bullet styles
\usepackage{graphicx}           % graphics
\usepackage{xcolor}             % color
\usepackage{hyperref}           % linkable references
\hypersetup{colorlinks=true,
            linkcolor=black,
            urlcolor=blue,
            citecolor=black
    }
  
%\usepackage{setspace}          % Together with \doublespacing below allows for
                                % doublespacing of the document

\usepackage{biblatex}
\addbibresource{references.bib}

\oddsidemargin=-0.5cm
\setlength{\textwidth}{7in}
\addtolength{\voffset}{-20pt}
\addtolength{\headsep}{25pt}
\renewcommand\thefootnote{[\arabic{footnote}]} % Bracketed footnotes

                                        % Macro definitions for the header
\newcommand{\myteam}{Miles Leonard,
                     Danny Nguyen}      % Name
\newcommand{\myclass}{DSA 5103}         % Class
\newcommand{\myassignment}{Project Report}  % Assignment designator

% Page headers automatically horizontal fill; Name ID Date Class Assignment Page.No.
\pagestyle{myheadings}
\markright{\myteam \hfill \today \hfill \myclass \hfill \myassignment \hfill }

\newcommand{\eqn}[0]{\begin{array}{rcl}}    % begin an aligned equation - allows
                                            % for aligning = or inequalities.
                                            % Always use with $$ $$
\newcommand{\eqnend}[0]{\end{array} }  	    % end the aligned equation

\newcommand{\qed}[0]{$\square$}     % make an unfilled square the
                                    % default for ending a proof

% Common number sets, must be in math mode
\newcommand{\Naturals}[0]{\mathbb{N}}   % Natural numbers
\newcommand{\Integers}[0]{\mathbb{Z}}   % Integers
\newcommand{\Rationals}[0]{\mathbb{Q}}  % Rationals
\newcommand{\Reals}[0]{\mathbb{R}}      % Reals
\newcommand{\Complex}[0]{\mathbb{C}}    % Complex

\newcommand{\heading}[1]{\noindent\large{\textbf{#1}} \normalsize}  % for sections
\newcommand{\problem}[1]{\noindent\paragraph{#1:}}      % for problem labels

%\doublespacing         % Together with the package setspace above allows for
                        % doublespacing of the document

\title{Modeling light pollution using geospatial, meteorologoical, energy usage, 
       and population data}
\author{\myteam\\
        Group 1\\
        Instructor: Dr. Charles Nicholson\\
        \myclass -- Fall 2025}
\date{\today}

\begin{document}

\maketitle
\pagebreak

\tableofcontents
\listoffigures
\listoftables
\pagebreak

\section{Executive Summary}

% 1 full page summarizing your work. Do not go over one page.
% – Concise problem statement
% – List of major concerns/assumptions (if any)
% – Summary of findings (include quantifiable information, e.g., CV R2 = 0.984)
% – Recommendations

\paragraph{Problem Statement} Light pollution is a night-time phenomenon
caused by the atmospheric scattering of excess light from things such as street 
lamps or other outdoor light fixtures that then bounces off of molecules in the air. 
Light pollution is commonly measured as night sky brightness (in magnitudes per square 
arcsecond), but is typically only directly measured at sites like observatories. 
Even then, data from these observatories is rarely accessible in real-time. The 
Globe at Night campaign created a monitoring network for observatories across the 
world to provide night sky brightness data for public use, with data going back 
as far as 2014. Light pollution, being an atmospheric phenomenon, is heavily 
impacted by weather conditions. The objective for this project is to predict light 
pollution (night sky brightness) at a given location using primarily meteorological, 
energy usage, and population data.

\paragraph{Major Concerns/Assumptions} TO-DO

\paragraph{Summary of Findings} TO-DO

\paragraph{Recommendations} TO-DO

\pagebreak

\section{Problem Background}

% – Problem description, context, background
% – Data description
% – Exploratory data analysis – the highlights; not the kitchen sink

\subsection{Problem Description}

\paragraph{}The use of artificial light at night (ALAN) in outdoor spaces—e.g., through the 
use of light fixtures such as street lamps or flood lights—causes what is known 
as light pollution. More specifically, improper and excessive use of ALAN, which 
has risen dramatically alongside rapid urbanization and industrialization around 
the globe, is the primary source of light pollution. The definition of light pollution 
varies across different agencies, but the common denominator amongst them all is 
that it is the culmination of the unwanted, adverse aspects of ALAN. There are 
three main components that make up light pollution: glare, light trespass, and 
skyglow; for the purposes of this project, however, the focus is solely on 
skyglow\cite{Welch}. Skyglow, the most broadly impactful aspect of light pollution, 
is the result of light being emitted upwards and then bouncing off of particles 
in the Earth's lower atmosphere, scattering. It obscures the natural night sky, 
preventing those in the city from being able to marvel at the stars. Those with 
exceptional vision and access to pristine night skies are able to see up to 7,000 
stars with the naked eye, however, that number plummets to fewer than 100 for the 
same observer in an urban area\cite{Mizon}.

\paragraph{}As it currently stands, light pollution-free skies are becoming more and more 
difficult to find. Over 99\% of US and European populations are affected by some 
level of light pollution, and a recent analysis found that skyglow is rapidly 
worsening at a rate of nearly 10\% per year\cite{Kyba}. Modeling the skyglow 
component of light pollution, typically referred to as night sky brightness, is 
a crucial aspect of many light pollution studies. Unfortunately, there is no 
definitive way to assess skyglow, and each method has its own pros and cons. 
Light pollution data for modeling often comes from a combination of ground-based 
measurements and satellite data. Ground-based measurements are most often taken 
using a sky quality meter (SQM), but have also been derived from citizen-science 
reports of naked-eye limiting magnitude (NELM). NELM is the magnitude (a 
logarithmic scale describing the apparent brightness of a star) of the faintest 
visible star to a naked-eye observer; the lower the magnitude, the brighter the 
object. The SQM is a handheld sensor that measures the brightness of the sky 
directly overhead in units of magnitudes per square arcsecond.

\paragraph{}Due to the complexity of light emission behavior and limitations in computing 
power, research into different modeling methods is still ongoing \cite{Kollath}. 
Most light pollution models are numerical models based on a series of functions 
that describe upward light flux and atmospheric scattering. Many more radiative 
transfer-based numerical models exist, but are often tailored towards modeling 
meteorological conditions and therefore not suitable for modeling light pollution 
specifically \cite{Kollath}. Recently, it was found that most, if not all light 
pollution models are significantly, systematically biased through implicitly 
assuming that all particles in the lower atmosphere are homogenous and spherical 
(Mie theory), despite being a safe assumption for daytime light models \cite{Kocifaj}. 
A popular global map known as the New world atlas of artificial night sky 
brightness \cite{Falchi} is another numerical model utilizing both ground-based 
and satellite measurements, but assumed a perfectly clear, calm, atmosphere over 
the entire globe. This limitation, albeit understandable, greatly impacts the 
accuracy of its output. Atmospheric conditions vary widely across regions. Not 
only that, but atmospheric conditions are always changing as well. Consecutive 
measurements of skyglow in the exact same location, even just hours apart, could 
produce different results.

\paragraph{}While acknowledging that the extent of light pollution is influenced by a 
multivariate set of economic indicators \cite{Olsen}, the current model 
specification is limited to assessing the predictive power of population metrics 
and energy consumption data.

\subsection{Data Description}

\paragraph{Night-Sky Brightness}Globe at Night (GaN) The GaN-MN project, an 
extension of the original Globe at Night project, is a global night sky brightness 
monitoring network using a commercially available meter (SQM-LE by Unihedron) for 
long-term monitoring of the light pollution conditions in different places around 
the world. \url{https://globeatnight.org/gan-mn/}

\paragraph{Meteorological Data} ERA5 provides hourly estimates for a large number 
of atmospheric, ocean-wave and land-surface quantities. An uncertainty estimate 
is sampled by an underlying 10-member ensemble at three-hourly intervals. 
Ensemble mean and spread have been pre-computed for convenience. Such uncertainty 
estimates are closely related to the information content of the available observing 
system which has evolved considerably over time. They also indicate flow-dependent 
sensitive areas. To facilitate many climate applications, monthly-mean averages 
have been pre-calculated too, though monthly means are not available for the 
ensemble mean and spread. \url{https://cds.climate.copernicus.eu/datasets/reanalysis-era5-pressure-levels?tab=overview}

\paragraph{Zip-Code Populations} A tabulation of the total population per zip code
per the 2020 U.S. Census. \url{https://data.census.gov/table}

\paragraph{Electric Retail Service Territories} A table of electric power retail 
service territories in the U.S. These are areas serviced by electric power utilities 
responsible for the retail sale of electric power to local customers, whether 
residential, industrial, or commercial. 
\url{https://ornl.opendatasoft.com/explore/dataset/electric-retail-service-territories/information/}

\paragraph{Population-Weighted Centroids for Zip Code Tabulation Areas (ZCTAs)} 
This dataset denotes ZIP Code centroid locations weighted by population. 
\url{https://hudgis-hud.opendata.arcgis.com/datasets/d032efff520b4bf0aa620a54a477c70e/about}

\subsection{Exploratory Data Analysis}

\paragraph{Meteorological Data} ERA5 data first look blurb here.

\paragraph{Electric-Service Provider Coverage} The service territories dataset is 
missing zip-code data for approximately 58\% of electrical service providers. This 
will limit the association of energy consumption estimates to only those zip-codes
that are in close proximity to providers with a known zip code.

\paragraph{Population} 4320 zip-codes were dropped between the intersection of
the centroids and population datasets. This is due to zip-codes that did not have
population data in the 2020 census data.

\section{Methodology}

The scope is geographically constrained to the 
contiguous United States, leveraging population data at the ZCTA resolution 
retrieved from the US Census Bureau. Energy consumption estimates will be derived 
by geospatial proximity analysis linking ZCTA-level data to the operational 
territories of adjacent utility providers. The energy consumption estimates sourced 
from the US Department of Energy Open Data Portal is missing location data for 
approximately 57\% of electrical service providers, resulting in a model training 
restriction to infer nearby population centers to be proportionate to the nearest 
electrical service providers.

% – Feature selection, engineering, missing value imputation, outlier processing, etc.
% – Modeling choices
% – State model validation approach (e.g., 5-fold CV) — note: Do not simply do a single holdout
% (80/20, 70/30, etc.) for anything unless you have a very good, defensible reason for doing so!

\section{Results}

% – Model results and performance summary
%   * for supervised learning models, this is pretty straightforward
%   * for unsupervised learning or other model types, you need to convince the reader of the validity
%     of your work
% – You should provide a detailed analysis of your best model (e.g., residual diagnostics, coefficients,
%   feature importance, hyperparameter tuning visualizations, confusion matrices, etc.)
% – Key findings of analysis! Tie this back to your models, assumptions, problem description, etc.
\subsection{Model Results and Performance}

\subsection{Key Findings}

\section{Conclusion}

% – Summary of problem, approach, findings – yes, this is a bit redundant, but this is your chance to
%   help the reader really connect the problem statement with your work.
% – Key insights, points out critical assumptions or limitations, etc.
% – Final recommendation and potential impact of the work – take into account the performance of
%   your models! For example, if your accuracy was 0.54, don’t say ”Company X can use this model
%   to improve their profits next quarter”, such a model is not very good, so they probably can’t do
%   that!

\paragraph{Summary}

\paragraph{Key Insights}

\paragraph{Final Recommendations and Impact}

\pagebreak

\printbibliography

\end{document}
